\section{Operator comparisons for a physical gap}\label{round26:app:operators}
\subsection{The direction of the comparison}
The first Theorem~11.1 of Faizal and Shabir~\cite{Round26FS26}, PDF
pages~130--131, uses a lower one-step comparison in (11.1).
Its proof then uses an upper Rayleigh comparison in (11.5).
Those two comparisons have opposite order. The following
example isolates the operator implication.

\begin{proposition}
There exist $a_k\downarrow0$ and bounded nonnegative
self-adjoint $H_k$ on one Hilbert space, all with the same
one-dimensional vacuum space, such that
\[
 T_k=e^{-a_kH_k},\qquad T_{k+1}\ge T_k,
\]
$H_0$ has gap one, and $H_k$ converges in operator norm to a
Hamiltonian with a unique vacuum and no positive spectral gap.
\end{proposition}
\begin{proof}
Take the orthonormal basis $\Omega,e_1,e_2,\ldots$ and define
\begin{gather*}
 a_k=2^{-k},\qquad H_k\Omega=0,\\
 H_ke_j=\max(a_k,1/j)e_j.
\end{gather*}
All positive eigenvalues of $H_0$ are one. The positive
eigenvalues of $H_k$ have infimum $a_k>0$, and the vacuum
is the span of $\Omega$. Both factors in
$a_k\max(a_k,1/j)$ decrease with $k$. Hence $T_{k+1}\ge T_k$
entry by entry in the common diagonal basis. The comparison
has identity compressions and zero defects.
Put $H\Omega=0$ and $He_j=j^{-1}e_j$. Then
$\|H_k-H\|\le a_k\to0$. The kernel of $H$ is still the
span of $\Omega$, but its positive eigenvalues $1/j$ tend
to zero. In fact $e^{-tH_k}\to e^{-tH}$ in operator norm
for every $t\ge0$, by the scalar mean-value bound on the
exponential. Thus even this strong convergence and exact
vacuum agreement do not repair the implication.
\end{proof}

The example tests the stated operator comparison. It does
not assert that these operators arise from the gauge-field
block map in~\cite{Round26FS26}.

\subsection{A comparison at equal physical times}
For a normalized vacuum $\Omega$, write
$P=|\Omega\rangle\langle\Omega|$ and $Q=I-P$.

\begin{lemma}\label{round26:lem:physical}
Let $H_k\ge0$ be self-adjoint on Hilbert spaces $\mathcal H_k$,
with vacuum projections $P_k$, and let $J_k:\mathcal H_{k+1}\to
\mathcal H_k$ be isometries satisfying $J_kP_{k+1}=P_kJ_k$.
Suppose $H_0\ge m_0Q_0$, $m_0>0$. Let $\delta_k\ge0$ satisfy
$\sum_k\delta_k<m_0$. If for some $a_{k+1}>0$ the quadratic
form bound
\begin{equation}\label{round26:eq:physical}
 Q_{k+1}e^{-a_{k+1}H_{k+1}}Q_{k+1}
 \le e^{a_{k+1}\delta_k}
       J_k^*Q_ke^{-a_{k+1}H_k}Q_kJ_k
\end{equation}
holds for every $k$, then
\[
 H_k\ge\left(m_0-\sum_{j<k}\delta_j\right)Q_k.
\]
\end{lemma}
\begin{proof}
The spectral theorem gives
$Q_ke^{-aH_k}Q_k\le e^{-am_k}Q_k$ whenever
$H_k\ge m_kQ_k$. Vacuum compatibility and the isometry
property give $J_k^*Q_kJ_k=Q_{k+1}$. Thus
\eqref{round26:eq:physical} bounds the left side by
$e^{-a_{k+1}(m_k-\delta_k)}Q_{k+1}$.
The scalar spectral theorem then implies
$H_{k+1}\ge(m_k-\delta_k)Q_{k+1}$. Induction proves the
claim. All these lower bounds are positive by the total-error
hypothesis.
\end{proof}

The right side compares semigroups at the same physical time.
Writing $T_k=e^{-a_kH_k}$, it involves
$T_k^{a_{k+1}/a_k}$, rather than $T_k$ without this time
conversion. An additive operator defect of size $\varepsilon_k$
generally produces a physical error of order
$\varepsilon_k/a_{k+1}$ when converted through the logarithm.
Summability of the dimensionless defects alone does not bound
this rescaled loss.

If all spaces are identified with one Hilbert space, the
vacuum projection is fixed, and $e^{-tH_k}$ converges strongly
to $e^{-tH}$ for each $t\ge0$, the uniform bound
$\|Qe^{-tH_k}Q\|\le e^{-tm_*}$ passes to the limit by
applying it to each vector. Therefore $H\ge m_*Q$ by the
spectral theorem. This last step requires a positive
$m_*=m_0-\sum_k\delta_k$, as supplied by the lemma.

\subsection{Selecting cross-cell covariance entries}
The argument after (5.23) of~\cite{Round26FS26}, PDF page~52,
uses positive semidefiniteness after selecting covariance
entries whose two supports belong to distinct cells. Such
entry selection requires a separate positivity proof.
For $0<r\le1$, the Gram matrix and its cross-cell mask are
\[
 C=\begin{pmatrix}1&r\\r&1\end{pmatrix},\qquad
 C_{\mathrm{cross}}=\begin{pmatrix}0&r\\r&0\end{pmatrix}.
\]
The first has eigenvalues $1+r,1-r$ and is positive
semidefinite. The second has eigenvalues $r,-r$.
It is not a principal submatrix of $C$. This algebraic
example concerns the entry-selection step, independently
of any subsequent field-theoretic estimate.

The missing input for the present Yang--Mills route remains a
bound for the actual Wilson measures that supplies a uniform
comparison of the required direction at equal physical times,
with a controlled total physical error and a nontrivial
continuum limit.

The finite-lattice identities are checked in \texttt{anc/check-slow-mode.py}. The coordinate heat-bath bound, the force overlap and the gap-transfer lemma have written proofs above; the checks are finite controls of their constants and normalization.
